\chapter{היבט עסקי ומסלול התרחבות}\label{ch:business}

\section{גודל השוק}
לפי \en{MarketsandMarkets}, שוק ה-\en{Deepfake AI}, שכולל גם פתרונות זיהוי ואימות, צפוי לצמוח מ-\num{857} מיליון דולר ב-\num{2025} ל-\num{7.27} מיליארד דולר ב-\num{2031}, בקצב שנתי של \num{42.8\%}~\cite{market2025}. זהו שוק רחב, שכולל גם זיהוי הונאות בבנקים ובביטוח. לכן \en{VeriLens} מגדירה את השוק שלה בצורה שמרנית (איור~\ref{fig:tam}):

\begin{figure}[!htb]
\begin{ltrdiagram}
% tam: TikZ / pgfplots source, built inside ltrdiagram (LTR)
\begin{tikzpicture}[font=\sffamily\small]
\fill[acc!12,draw=acc] (0,0) circle (2.6);
\fill[acc!30,draw=acc] (0,-0.75) circle (1.8);
\fill[acc2!70,draw=acc2] (0,-1.55) circle (0.95);
\node[align=center] at (0,1.85) {\en{\textbf{TAM}}\\\he{שוק ה-\en{Deepfake AI}}\\\he{7.27 מיליארד \$ עד 2031}};
\node[align=center] at (0,0.3) {\en{\textbf{SAM}}\\\he{אימות מדיה לחדשות,}\\\he{פלטפורמות ופרסום}};
\node[align=center,text=white] at (0,-1.55) {\en{\textbf{SOM}}\\\he{1,000 מערכות}};
\node[align=left,anchor=west,font=\sffamily\scriptsize,text width=5.6cm] at (3.2,0.6) {\he{\textbf{\en{SAM}}: אומדן עבודה של כ-10\% מה-\en{TAM}. יאומת בדוח ההתקדמות מול נתוני לקוחות.}};
\node[align=left,anchor=west,font=\sffamily\scriptsize,text width=5.6cm] at (3.2,-1.4) {\he{\textbf{\en{SOM}}: 1,000 מערכות $\times$ מנוי ממוצע של 3,500 ₪ בחודש $\approx$ 42 מיליון ₪ בשנה (חודש 36).}};
\end{tikzpicture}

\end{ltrdiagram}
\caption{גודל השוק: גם נתח קטן מהשוק מספיק להכנסה של \num{42} מיליון ₪ בשנה}\label{fig:tam}
\end{figure}

\section{מודל הכנסות}

המודל נשען על שלושה מקורות הכנסה, שמפורטים בטבלה~\ref{tab:revenue}.

\begin{table}[!htb]\centering
\caption{שלושה מקורות הכנסה}\label{tab:revenue}
\begin{tabularx}{\linewidth}{@{}l Y l@{}}
\toprule
\textbf{מקור} & \textbf{לקוח ומה הוא מקבל} & \textbf{תמחור} \\
\midrule
מנוי ענן & מערכות חדשות: פלטפורמה, מסוף ולוח בקרה & \num{1,500} עד \num{15,000} ₪ בחודש לפי נפח \\
הסמכת \en{TRUTH-IMG} & מערכות חדשות: ביקורת שנתית ותג ״מערכת מאומתת״ & \num{25,000} ₪ בשנה \\
ממשק ״ציון שלמות״ & רשתות פרסום ומפרסמים: \en{Brand Safety} & לפי אלף שאילתות \\
\bottomrule
\end{tabularx}
\end{table}

התמחור נבדק מול המתחרה הישיר: \en{ImageWhisperer} גובה \num{500} עד \num{900} יורו בחודש עבור \num{1,500} עד \num{6,000} בדיקות~\cite{imagewhisperer}. המנוי החודשי שהנחנו (\num{3,500} ₪ בממוצע) הוא בסדר גודל דומה, אבל הוא כולל נפח גדול בהרבה, כך שהמחיר לתמונה אצלנו נמוך יותר. זה יתרון מכירה, אבל גם סיכון: אם הלקוחות ישתמשו בנפח קטן, ההכנסה לתמונה תעלה והמודל עדיין יחזיק.

מקור ההכנסה השלישי הוא המנוף האתי והכלכלי של הפרויקט. אם מפרסמים יעדיפו אתרים עם ציון שלמות גבוה, ההיגיון הכלכלי של הלולאה באיור~\ref{fig:loop} מתהפך: תמונה מזויפת תוריד את הכנסות האתר במקום להעלות אותן.

\section{מי עוד משלם: שווקים חלופיים}
למערכות חדשות יש צורך אמיתי אבל תקציבים קטנים, והכלי החינמי לעיתונאים \en{TrueMedia} נסגר אחרי פחות משנה~\cite{truemedia2025}. לכן הטכנולוגיה מתוכננת כך שתשרת גם לקוחות שתמונה מזויפת עולה להם כסף באופן ישיר. הטבלה היא הערכה של הצוות, שתיבדק בשיחות עם לקוחות בחודשים הראשונים.

\begin{table}[!htb]\centering
\caption{שווקים חלופיים לאותה טכנולוגיה (הערכה, לבדיקה בשיחות עם לקוחות)}\label{tab:altmarkets}
\small
\begin{tabularx}{\linewidth}{@{}l Y l@{}}
\toprule
\colhead{לקוח} & \colhead{למה הוא צריך} & \colhead{חוזק כלקוח משלם} \\
\midrule
חברות ביטוח & תביעות עם תמונות נזק מזויפות; יומן ההחלטות משמש ראיה & חזק \\
אתרי מסחר ושליחויות & תמונה מזויפת של מוצר ״פגום״ לקבלת החזר & חזק \\
סוכנויות צילום וידיעות & כל הערך שלהן הוא שהתמונה אמיתית & בינוני \\
מערכות חדשות & צורך אמיתי, תקציבים קטנים & חלש עד בינוני \\
\bottomrule
\end{tabularx}
\end{table}

\section{השוק הישראלי: למה מתחילים כאן}\label{sec:israel}
בישראל הבעיה כבר לא תיאורטית, ושני אירועים מהשנתיים האחרונות מראים את זה. במלחמה עם איראן ביוני \num{2025} הופצו ברשתות סרטונים ותמונות שנוצרו בבינה מלאכותית של ״פגיעות״ בתל אביב ובנתב״ג, ו-\en{NewsGuard} מצאה \num{51} אתרים שהפיצו טענות שקריות, בהן תמונות הרס מזויפות~\cite{afp2025}. בבחירות לכנסת ה-\num{26} (\num{27} באוקטובר \num{2026}) קבעה ועדת הבחירות המרכזית, בראשות השופט נעם סולברג, חובת סימון לתעמולה שנוצרה או שונתה באופן דיגיטלי: הודעה בולטת או אחד משני סמלים אחידים, ״נוצר דיגיטלית״ ו״שונה דיגיטלית״~\cite{pearlcohen2026,toi2026}. במהלך הקמפיין אף הורתה הוועדה להסיר תמונה מעובדת שפרסמה מפלגה~\cite{toi2026}.

המשמעות לפרויקט: הציבור הישראלי כבר מכיר תווית ״נוצר דיגיטלית״ מהבחירות. התוויות של \en{VeriLens} יתואמו לסמלים של ועדת הבחירות, כדי שהקורא יראה סמל אחד ומוכר, ולא סמל חדש בכל אתר.

\textbf{המתחרים המקומיים.} \en{Cyabra} הישראלית השיקה ביולי \num{2025} כלי לזיהוי תמונות וסרטונים מזויפים, אבל הוא מכוון למותגים ולממשלות ולא למערכות חדשות~\cite{cyabra2025}. \en{Clarity} הישראלית, שהתמחתה בזיהוי זיופים, נרכשה באוגוסט \num{2026} והתמקדה מאז באימות זהות של עובדים~\cite{clarity2026}. כלומר, אין היום ספק ישראלי שמוכר אימות תמונות למערכות חדשות, בעברית ובערבית.

\begin{table}[!htb]\centering
\caption{השוק הישראלי: מי הלקוחות, לפי סדר הפנייה}\label{tab:israel}
\small
\begin{tabularx}{\linewidth}{@{}l Y Y@{}}
\toprule
\colhead{לקוח} & \colhead{למה הוא צריך} & \colhead{איך פונים} \\
\midrule
אתרי חדשות ארציים & מלחמות ובחירות מביאות גל של תמונות מזויפות, והטעות מתפרסמת מול מיליוני קוראים & פיילוט חינמי לדסק אחד, דרך שלוש מערכות הפיילוט (\req{OP-04}) \\
תקשורת בערבית בישראל & אותה בעיה, ואין היום כלי בערבית שמותאם לה & ממשק ותוויות בערבית מהיום הראשון \\
גופי בדיקת עובדות & בודקים תמונות ביד, תמונה אחר תמונה & רישיון מוזל, בתמורה למשוב ולשיתוף מאגר זיופים \\
יחידות דוברות ויחסי ציבור & תמונה שמפיצים בשם הארגון חייבת להיות אמיתית & מנוי לפי נפח, בלי שילוב במערכת ניהול תוכן \\
חברות ביטוח & תביעות עם תמונות נזק מזויפות (טבלה~\ref{tab:altmarkets}) & אותה טכנולוגיה, ממשק נפרד, אחרי הפיילוט בחדשות \\
\bottomrule
\end{tabularx}
\end{table}

השוק הישראלי קטן מדי כדי לבנות עליו חברה לבדו, ולכן הוא שלב ראשון ולא היעד (איור~\ref{fig:roadmap}). הוא מתאים לשלב הראשון משלוש סיבות: הצורך כבר הוכח במלחמה ובבחירות, אנחנו מכירים את השפה ואת השוק, והמערכות קטנות מספיק כדי שפיילוט יגיע מהר לעורך הראשי. מה שעוד צריך להוכיח הוא שמישהו ישלם, וזו מטרת השיחות ב-\req{OP-01}.

\section{מפת דרכים לחמש שנים}

ההתרחבות מותנית: עוברים לשלב הבא רק אחרי שעוברים את השער של השלב הקודם, כפי שמראה איור~\ref{fig:roadmap}.

\begin{figure}[!htb]
\begin{ltrdiagram}
% roadmap: TikZ / pgfplots source, built inside ltrdiagram (LTR)
\begin{tikzpicture}[font=\sffamily\scriptsize,
  ph/.style={draw=acc,fill=soft,rounded corners=4pt,text width=2.7cm,minimum height=2.3cm,align=center,inner sep=4pt},
  gt/.style={diamond,fill=acc2,inner sep=2.2pt}]
\draw[very thick,ink,-{Stealth}] (-0.2,0)--(15.4,0);
\foreach \x/\t/\b/\g in {
 1.4/{\he{\textbf{שנה 1 עד 1.5}}}/{\he{פיתוח ופיילוט}\\\he{5 מערכות בישראל}\\\he{2 מיליון תמונות/חודש}}/{\he{קבלה מבצעית}},
 4.5/{\he{\textbf{שנה 2}}}/{\he{אירופה, סעיף 50}\\\he{100 מערכות}\\\he{10 מיליון/חודש}}/{\he{100 לקוחות משלמים}},
 7.6/{\he{\textbf{שנה 3}}}/{\he{גלובלי, 30 מדינות}\\\he{1,000 מערכות}\\\he{50 מיליון/חודש}}/{\he{איזון תפעולי}},
 10.7/{\he{\textbf{שנה 4}}}/{\he{וידאו ואודיו}\\\he{שותפות עם רשת פרסום}}/{\he{גלאי וידאו בתקן}},
 13.8/{\he{\textbf{שנה 5}}}/{\he{\en{TRUTH-IMG} כתקן}\\\he{תעשייתי מוכר}}/{\he{הכרה רשמית}}}{
 \node[ph,anchor=south] at (\x,0.35) {\t\\[2pt]\b};
 \node[gt] at (\x,0) {};
 \node[font=\sffamily\scriptsize,text=acc2,align=center] at (\x,-0.45) {\g};}
\end{tikzpicture}

\end{ltrdiagram}
\caption{מפת דרכים: מתרחבים רק אחרי שעוברים את שער השלב הקודם (המעוינים)}\label{fig:roadmap}
\end{figure}

כל שלב נפתח רק אחרי שעבר את השער שלו. כך החזון הגדול נשאר מחויב למספרים: לא מתרחבים לאירופה בלי קבלה מבצעית, ולא נכנסים לווידאו לפני שהתמונות רווחיות.
